\documentclass[11pt,a4paper]{article}

\usepackage{fontspec}
\usepackage[ngerman]{babel}
\usepackage{geometry}
\geometry{margin=2.6cm}
\usepackage{amsmath,amssymb,mathtools}
\usepackage{enumitem}
\usepackage{xcolor}
\usepackage[unicode,colorlinks=true,linkcolor=blue!40!black,urlcolor=blue!40!black]{hyperref}

\title{Die neue Physik}
\author{Thomas Hoffbauer}
\date{}

\begin{document}
\maketitle

\begin{center}
\large\itshape
Ein mathematisches Gedankenexperiment über Primzahlen, Oktonionen, \(E_8\)\\
und die Frage, wie weit sich Naturgesetze aus reiner Struktur ableiten lassen
\end{center}
\vspace{1.25em}

Es gibt in der modernen theoretischen Physik zwei grundsätzlich verschiedene Wege, nach einer tieferen Ordnung der Natur zu suchen. Der eine beginnt bei Experimenten und versucht, aus Messdaten die zugrunde liegenden Gesetze zu rekonstruieren. Der andere beginnt bei der Mathematik selbst: bei Symmetrien, Zahlentheorie, Gittern und Operatoralgebren – und fragt, ob bestimmte physikalische Strukturen vielleicht nicht bloß beschrieben, sondern durch Mathematik nahezu erzwungen werden.


Ein ungewöhnliches Forschungsprogramm, das wir hier als EABC-Modell bezeichnen, verfolgt konsequent den zweiten Weg. Sein Ausgangspunkt ist keine neue Elementarteilchenhypothese. Es beginnt vielmehr mit einigen der außergewöhnlichsten Objekte der reinen Mathematik: dem \(E_8\)-Gitter, den Oktonionen, der Möbiusfunktion, der Primfaktorzerlegung ganzer Zahlen und den Clifford-Algebren.


Das Resultat ist bislang keine neue fundamentale Theorie der Natur. Aber es entsteht ein erstaunlich dichtes Netz exakter mathematischer Beziehungen. Manche davon sind klassische Sätze. Einige konnten in einem Lean-4-Theorem-Prover formal nachvollzogen werden. Andere bleiben Modellhypothesen. Und gerade diese Trennung ist entscheidend.


Die zentrale Frage lautet:


Wie viel von der formalen Architektur moderner Physik lässt sich allein aus Arithmetik und Geometrie rekonstruieren – bevor überhaupt physikalische Postulate eingeführt werden?





\section{Eine Welt aus 240 Berührungen}


Das geometrische Zentrum des Programms ist das \(E_8\)-Gitter.


Es lebt im acht-dimensionalen euklidischen Raum und besitzt 240 Vektoren minimaler Länge. Für jeden dieser Wurzelvektoren v gilt


\[
\|v\|^2=2.
\]


Legt man um jeden Gitterpunkt eine Kugel vom Radius


\[
r=\frac{1}{\sqrt{2}},
\]


so berühren genau 240 Nachbarkugeln die Kugel im Ursprung. Ihre Mittelpunkte liegen im Abstand \(\sqrt{2}\), die tatsächlichen Kontaktpunkte bei


\[
\frac12v.
\]


Damit besitzt die \(E_8\)-Packung Kusszahl


\[
K(8)=240.
\]


Maryna Viazovska zeigte 2016, dass diese Konfiguration nicht nur hochsymmetrisch ist. Die \(E_8\)-Packung ist tatsächlich die dichteste Kugelpackung im acht-dimensionalen Raum.


Ihre Packungsdichte beträgt


\[
\Delta_8=\frac{\pi^4}{384}.
\]


Die 240 Minimalvektoren bilden zugleich das Wurzelsystem der außergewöhnlichen Lie-Algebra \(E_8\).


Schon hier begegnen sich drei ansonsten weit voneinander entfernte Gebiete:


\[
\text{Kugelpackung}
\longleftrightarrow
\text{Gittertheorie}
\longleftrightarrow
\text{Lie-Algebra}.
\]





\section{Die 240 Wurzeln sind hochgradig strukturiert}


In einer Standardkoordinatisierung zerfallen die \(E_8\)-Wurzeln in zwei Klassen.


112 besitzen die Gestalt


\[
(\pm1,\pm1,0,0,0,0,0,0),
\]


wobei die beiden von null verschiedenen Positionen beliebig gewählt werden dürfen.


Denn


\[
\binom82\,2^2=28\cdot4=112.
\]


Die übrigen 128 haben die Form


\[
\frac12(\pm1,\pm1,\pm1,\pm1,\pm1,\pm1,\pm1,\pm1),
\]


mit einer geraden Zahl negativer Vorzeichen.


Damit gilt


\[
240=112+128.
\]


Diese Zerlegung ist nicht bloße Buchhaltung. Sie verbindet \(E_8\) mit Spinordarstellungen, Binärcodes und den speziellen Symmetrien acht-dimensionaler Räume.


Eine weitere bemerkenswerte Struktur ist eine Zerlegung in orthogonale Bezugssysteme. Der 112er-Sektor kann über eine 1-Faktorisierung des vollständigen Graphen \(K_8\) in sieben Gruppen zu je 16 Wurzeln organisiert werden.


Ein typisches System lautet


\[
\{\pm e_0\pm e_1,\;
\pm e_2\pm e_3,\;
\pm e_4\pm e_5,\;
\pm e_6\pm e_7\}.
\]


Die vier Achsen sind paarweise orthogonal.


Die sieben solchen Systeme lassen sich elegant mit der Fano-Ebene verbinden – jener endlichen projektiven Ebene mit sieben Punkten und sieben Geraden, die zugleich die Multiplikationsstruktur der Oktonionen kodiert.


Der 128er-Sektor wiederum besitzt eine natürliche Beziehung zu Hadamard-Strukturen und binären Codes.


Das geometrische Bild beginnt sich zu verdichten:


\[
E_8
\quad\leftrightarrow\quad
\text{Fano-Geometrie}
\quad\leftrightarrow\quad
\text{Codes}
\quad\leftrightarrow\quad
\text{Spinoren}.
\]





\section{\texorpdfstring{Von drei Dimensionen zu \(E_8\)}{Von drei Dimensionen zu E8}}


Eine besonders eindrucksvolle Brücke stammt aus Arbeiten von Pierre-Philippe Dechant.


Man beginnt mit dem ikosaedrischen Wurzelsystem \(H_3\) im dreidimensionalen Raum. Seine Wurzeln werden in der Clifford-Algebra \(\mathrm{Cl}(3)\) multipliziert.


Die geradgradigen Produkte bilden Spinoren. Im Fall der Ikosaedersymmetrie entstehen 120 solcher Spinoren – die binäre Ikosaedergruppe.


Der gerade Teil von \(\mathrm{Cl}(3)\) ist vierdimensional:


\[
\mathrm{Cl}^{+}(3)\cong\mathbb H.
\]


Interpretiert man die 120 Spinoren als Vektoren in diesem vierdimensionalen Raum, erhält man das außergewöhnliche Wurzelsystem \(H_4\), dessen 120 Wurzeln zugleich die Ecken des 600-Zells bilden.


Schematisch:


\[
H_3
\longrightarrow
\text{Spinoren in }\mathrm{Cl}^{+}(3)
\longrightarrow
H_4.
\]


Nimmt man anschließend die volle acht-dimensionale Clifford-Algebra \(\mathrm{Cl}(3)\), erhält man 240 achtkomponentige Clifford-Objekte. Mit einem speziell reduzierten Skalarprodukt über den goldenen Schnitt


\[
\tau=\frac{1+\sqrt{5}}{2}
\]


entsteht daraus ein Wurzelsystem vom Typ \(E_8\).


Damit existiert tatsächlich eine mathematisch harte Route


\[
3\mathrm{D}
\longrightarrow
4\mathrm{D}
\longrightarrow
8\mathrm{D},
\]


die nicht auf der numerischen Beobachtung


\[
24\times10=240
\]


beruht.


Die 240 \(E_8\)-Wurzeln besitzen also eine explizite dreidimensional-spinorielle Realisierung.





\section{Was noch fehlt: die explizite Isometrie}


Die Dechant-Konstruktion erzeugt ein \(E_8\)-Wurzelsystem. In einem formalen Projekt besitzt man aber häufig bereits eine andere Koordinatendarstellung von \(E_8\).


Dann stellt sich eine konkrete Frage:


Gibt es eine explizite lineare Abbildung


\[
T:\Phi_{\mathrm{Dechant}}
\longrightarrow
\Phi_{E_8}^{\mathrm{repo}},
\]


die beide Darstellungen punktweise miteinander identifiziert?


Man wählt acht einfache Wurzeln und schreibt sie als Spalten zweier Matrizen


\[
B_D,\qquad B_R.
\]


Erfüllen beide dieselbe Gram-Gleichung


\[
B_D^TB_D=B_R^TB_R=C_{E_8},
\]


so definiert


\[
T=B_RB_D^{-1}
\]


den natürlichen Basiswechsel.


Dann gilt


\[
T\alpha_i^D=\alpha_i^R.
\]


Und aus der Gram-Gleichheit folgt exakt


\[
T^TT=I.
\]


Ein numerischer Test ergab eine Abweichung


\[
\|T^TT-I\|\approx7{,}1\times10^{-15},
\]


also Maschinengenauigkeit.


Aber hier wird die methodische Trennung wichtig:


Numerisch extrem genau ist nicht dasselbe wie formal bewiesen.


Der mathematische Weg zur vollständigen Identifikation führt über Weyl-Reflexionen. Für eine Isometrie gilt


\[
T\,s_\alpha=s_{T\alpha}\,T.
\]


Werden die einfachen Reflexionen korrekt transportiert, folgt dasselbe für jedes Weyl-Wort. Damit lässt sich schließlich zeigen


\[
T(\Phi_D)=\Phi_R.
\]


Ein direkter 240-Punkte-Abgleich wäre dann lediglich das abschließende Rechenaudit.





\section{\texorpdfstring{\(E_8\) besitzt eine arithmetische Schalenstruktur}{E8 besitzt eine arithmetische Schalenstruktur}}


Die 240 Wurzeln sind nur die erste Schale des Gitters.


Die Theta-Reihe von \(E_8\) besitzt die bemerkenswerte Identität


\[
\Theta_{E_8}(\tau)=E_4(\tau),
\]


wobei \(E_4\) die Eisensteinreihe vom Gewicht vier ist.


Damit gilt


\[
\Theta_{E_8}
=
\]

1+

\[
240\sum_{n\ge1}\sigma_3(n)q^n,
\]


und folglich


\[
\boxed{
N_{E_8}(n)=240\,\sigma_3(n)
}
\]


für die Anzahl der Gitterpunkte mit Normquadrat 2n.


Hier taucht plötzlich Zahlentheorie mitten in der acht-dimensionalen Geometrie auf.


Die Funktion


\[
\sigma_3(n)=\sum_{d\mid n}d^3
\]


ist multiplikativ. Die Besetzung jeder geometrischen Schale kennt also die Primfaktorzerlegung ihres Index.


Das ist kein EABC-Postulat. Es ist klassische Mathematik.





\section{Warum Viazovskas Beweis so außergewöhnlich ist}


Die Optimalität von \(E_8\) lässt sich durch das Cohn-Elkies-Verfahren in ein Optimierungsproblem für Funktionen übersetzen.


Man sucht eine radialsymmetrische Schwartz-Funktion f mit


\[
f(x)\le0
\qquad
\text{für }\|x\|\ge\sqrt{2}
\]


und


\[
\widehat f(y)\ge0.
\]


Gelingt es zusätzlich, die Schranke exakt bei der \(E_8\)-Dichte zu sättigen, ist die Kugelpackung optimal.


Viazovskas „magische Funktion“ besitzt Nullstellen genau an den \(E_8\)-Schalen.


Der charakteristische Faktor


\[
\sin^2\left(\frac{\pi r^2}{2}\right)
\]


erzeugt Nullstellen bei


\[
r=\sqrt{2n}.
\]


Die eigentliche Konstruktion beruht auf Modul- und Quasimodulformen. Dadurch begegnen sich erneut Geometrie, Fourieranalysis und arithmetische Funktionen.


Die Dichte folgt schließlich als


\[
\Delta_8
=
\frac{\pi^4}{384}.
\]


In 24 Dimensionen geschieht ein verwandtes Wunder für das Leech-Gitter, dessen erste Schale 196560 Minimalvektoren besitzt.


Dimension acht und Dimension 24 sind damit außergewöhnliche Fixpunkte in der Theorie optimaler Packungen.





\section{\texorpdfstring{Vom \(E_8\)-Spektrum zur Cyclotomie}{Vom E8-Spektrum zur Cyclotomie}}


Das \(E_8\)-Wurzelsystem besitzt Coxeter-Zahl


\[
h=30.
\]


Seine Exponenten sind


1,7,11,13,17,19,23,29.


Bemerkenswerterweise ist dies genau


\[
(\mathbb Z/30\mathbb Z)^\times.
\]


Die Eigenwerte eines Coxeter-Elements liegen daher im 30. Kreisteilungskörper


\[
\mathbb Q(\zeta_{30}).
\]


Seine Galoisgruppe ist


\[
\operatorname{Gal}
\bigl(
\mathbb Q(\zeta_{30})/\mathbb Q
\bigr)
\cong
(\mathbb Z/30\mathbb Z)^\times.
\]


Nun betrachtet man die Quadratuntergruppe


\[
G^2=\{1,19\}.
\]


Der Quotient besitzt vier Elemente:


\[
G/G^2\cong V_4.
\]


Damit taucht die Kleinsche Vierergruppe innerhalb der cyclotomischen Struktur des \(E_8\)-Coxeter-Spektrums auf – diesmal nicht als numerisches Muster, sondern als echter Gruppenquotient.


Was bislang nicht gezeigt ist: ob genau dieses \(V_4\) geometrisch durch den Normalisator eines Coxeter-Zyklus in der Weylgruppe realisiert wird.


Das entsprechende Gate bleibt offen.





\section{Primzahlen als Besetzungsplätze}


Die vielleicht überraschendste Entwicklung beginnt nicht bei \(E_8\), sondern bei der elementaren Primfaktorzerlegung.


Jede quadratfreie Zahl


\[
n=p_{i_1}p_{i_2}\cdots p_{i_k}
\]


entspricht eindeutig einer endlichen Menge von Primzahlen.


Damit besitzt


\[
\ell^2(\mathbb N_{\mathrm{sf}})
\]


dieselbe Besetzungskombinatorik wie ein fermionischer Fockraum.


Jeder Primplatz ist entweder unbesetzt oder einmal besetzt:


\[
N_p\in\{0,1\}.
\]


Mit einer geordneten Primzahlbasis können Erzeugungsoperatoren durch ein Jordan-Wigner-artiges Vorzeichen definiert werden:


\[
a_{p_k}^\dagger|S\rangle
=
(-1)^{\#\{p_j\in S:j<k\}}
|S\cup\{p_k\}\rangle.
\]


Dann entstehen die kanonischen Antikommutationsrelationen


\[
(a_p^\dagger)^2=0,
\]


\[
\{a_p^\dagger,a_q^\dagger\}=0,
\]


\[
\{a_p,a_q^\dagger\}
=
\delta_{pq}.
\]


In dem hier diskutierten Formalisierungsprojekt wurden solche Relationen in Lean 4 implementiert.


Besonders bemerkenswert ist die Fermionparität.


Für quadratfreie Zahlen gilt


\[
\mu(n)=(-1)^{\omega(n)}.
\]


Damit lässt sich schreiben


\[
(-1)^F|n\rangle
=
\mu(n)|n\rangle.
\]


Die Möbiusfunktion übernimmt innerhalb dieses arithmetisch konstruierten Fockraums exakt die Rolle der Fermionparität.


Das ersetzt nicht das physikalische Spin-Statistik-Theorem.


Aber es zeigt etwas mathematisch Eigenständiges:


Die multiplikative Struktur quadratfreier Zahlen besitzt eine natürliche CAR-Darstellung.





\section{Woher könnte der bosonische Sektor kommen?}


Jeder Primexponent lässt sich eindeutig schreiben als


\[
e_p=f_p+2k_p,
\]


mit


\[
f_p\in\{0,1\},
\qquad
k_p\in\mathbb N.
\]


Damit zerfällt jede positive ganze Zahl eindeutig in


\[
n=n_{\mathrm{sf}}m^2.
\]


Der erste Faktor besitzt fermionische Besetzungskombinatorik.


Der zweite Faktor besitzt unbeschränkte Besetzungszahlen


\[
k_p=0,1,2,\ldots,
\]


also die Kombinatorik bosonischer Moden.


Schematisch:


\[
\ell^2(\mathbb N)
\cong
\ell^2(\mathbb N_{\mathrm{sf}})
\widehat\otimes
\mathcal H_{\mathrm{bos}}.
\]


Damit entsteht ein natürlicher Kandidat für eine arithmetische Bose-Fermi-Zerlegung.


Aber auch hier gilt eine harte Grenze:


Unbeschränkte Besetzungszahlen sind noch keine bosonische Quantenfeldtheorie.


Dafür müssen Operatoren konstruiert werden, die tatsächlich


\[
[b_p,b_q^\dagger]
=
\delta_{pq}
\]


erfüllen.


Erst wenn CAR und CCR gemeinsam vorliegen, kann sinnvoll von einer arithmetischen Superalgebra gesprochen werden.





\section{Eine mögliche arithmetische Supersymmetrie}


Nimmt man einen Fermionoperator \(a_p\) und einen Bosonoperator \(b_p\), ergibt sich die klassische supersymmetrische Konstruktion


\[
Q_p=a_pb_p^\dagger,
\qquad
Q_p^\dagger=a_p^\dagger b_p.
\]


Wegen


\[
a_p^2=(a_p^\dagger)^2=0
\]


folgt


\[
Q_p^2=(Q_p^\dagger)^2=0.
\]


Der Antikommutator ergibt


\[
\{Q_p,Q_p^\dagger\}
=
N_{b,p}+N_{f,p}.
\]


Damit entsteht ein natürlicher supersymmetrischer Hamiltonoperator.


Interessanterweise unterscheidet er sich vom arithmetischen Primexponenten


\[
e_p=2N_{b,p}+N_{f,p}.
\]


Das ist kein Schönheitsfehler. Es zeigt vielmehr, dass zwei verschiedene natürliche Gradierungen existieren:


\[
H_p^{\mathrm{SUSY}}
=
N_{b,p}+N_{f,p}
\]


und


\[
E_p^{\mathrm{arith}}
=
2N_{b,p}+N_{f,p}.
\]


Die Mathematik zwingt also nicht sofort eine physikalische Energieinterpretation auf.





\section{Oktonionen: sieben imaginäre Richtungen}


Ein weiterer Strang beginnt bei den Oktonionen


\[
\mathbb O
=
\mathbb R
\oplus
\operatorname{Im}(\mathbb O).
\]


Ihre Automorphismengruppe ist


\[
G_2=\operatorname{Aut}(\mathbb O).
\]


Sie besitzt Dimension 14.


Fixiert man eine normierte imaginäre Oktonioneneinheit u, so ist ihr Stabilisator


\[
\operatorname{Stab}_{G_2}(u)\cong SU(3).
\]


Daraus folgt die bekannte Darstellung


\[
\mathbf7
\longrightarrow
\mathbf1\oplus\mathbf3\oplus\bar{\mathbf3}.
\]


Das ist eine echte mathematische Beziehung zwischen Oktonionen und \(\mathrm{SU}(3)\).


Sie bedeutet jedoch nicht automatisch, dass die drei komplexen Richtungen physikalische Quarkfarben sind.


Gerade hier ist die Versuchung zur Überinterpretation besonders groß.


Man kann drei komplexe Paare definieren und ihnen die Namen Rot, Grün und Blau geben. Man kann anschließend diagonale Operatoren mit Eigenwerten


\[
\frac12,\quad-\frac12,\quad0
\]


oder


\[
\frac{1}{2\sqrt{3}},
\quad
\frac{1}{2\sqrt{3}},
\quad
-\frac{1}{\sqrt{3}}
\]


konstruieren.


Damit lassen sich mathematisch dieselben Eigenwertmuster wie bei den Cartan-Generatoren von \(\mathfrak{su}(3)\) erzeugen.


Aber ein Lean-Satz, der drei unterschiedliche Datenstrukturen Red, Green und Blue auseinanderhält, beweist noch keine QCD.


Formalisiert werden muss vielmehr die tatsächliche Lie-Algebra-Wirkung:


\[
[T_a,T_b]
=
if_{abc}T_c.
\]


Erst dann existiert wirklich eine \(\mathfrak{su}(3)\)-Darstellung.





\section{Die entscheidende Grenze}


Viele der diskutierten Konstruktionen besitzen eine verblüffende formale Ähnlichkeit mit moderner Physik:


\begin{itemize}[leftmargin=*]
\item CAR aus quadratfreien Zahlen,
\item bosonenartige Besetzungszahlen aus Primzahlquadraten,
\item \(\mathrm{SU}(3)\) als Stabilisator in \(G_2\),
\item \(E_8\)-Spinoren aus dreidimensionaler Ikosaedergeometrie,
\item cyclotomische \(V_4\)-Strukturen,
\item Theta-Reihen und Zeta-Funktionen,
\item natürliche Spektren über \(\log p\).
\end{itemize}

Aber Ähnlichkeit ist keine Identität.


Die zentrale epistemische Firewall des Projekts lautet deshalb:


\[
\boxed{
\text{Mathematische Struktur}
\neq
\text{physikalische Interpretation}.
}
\]


Ein formal bewiesener Operator


\[
Y=\frac{2}{\sqrt{3}}T_8
\]


ist zunächst nur ein Operator.


Er wird nicht dadurch zur schwachen Hyperladung, dass seine Eigenwerte bekannte rationale Zahlen reproduzieren.


Ebenso ist


\[
G_2/SU(3)\cong S^6
\]


ein tiefer Satz der Differentialgeometrie.


Er beweist aber noch nicht, dass \(G_2\) die Eichgruppe einer fundamentalen Naturtheorie ist.





\section{Warum die formale Verifikation trotzdem wichtig ist}


Gerade deshalb ist der Einsatz eines Theorem-Provers wie Lean interessant.


In konventionellen spekulativen Theorien können mathematische Identitäten, numerische Zufälle und physikalische Interpretation unbemerkt ineinanderfließen.


Ein Proof Assistant zwingt zur Trennung.


Ein Satz ist entweder aus Definitionen und vorherigen Theoremen hergeleitet – oder nicht.


Man kann deshalb für jedes Modul ein Ledger führen:


MATHEMATICAL THEOREM

EXACT EXTERNAL CHECK

NUMERICAL PASS

LEAN KERNEL PROOF

MODEL HYPOTHESIS

PHYSICAL INTERPRETATION


Diese Kategorien sind nicht austauschbar.


Ein numerischer Test mit


\[
10^{-15}
\]


Fehler ist hervorragend.


Aber er ist kein exakter Beweis.


Eine Formel, die per Definition gesetzt wurde, ist exakt wahr.


Aber sie ist keine Vorhersage.


Und ein Lean-Theorem kann vollständig korrekt sein, obwohl die physikalische Interpretation des Theorems falsch ist.





\section{Die vielleicht interessanteste Frage}


Was bleibt also von der „neuen Physik“, wenn man alle überzogenen Behauptungen entfernt?


Erstaunlich viel.


Es bleibt ein Forschungsprogramm, das drei außergewöhnliche mathematische Welten zusammenführt:


\[
\boxed{
\text{arithmetische Fockstrukturen}
}
\]


\[
\boxed{
\text{exceptionale Geometrie }
G_2,E_8
}
\]


\[
\boxed{
\text{modulare Analysis und optimale Gitter}
}
\]


und untersucht, welche gemeinsamen Operatorstrukturen daraus hervorgehen.


Vielleicht entsteht daraus niemals eine neue fundamentale Physik.


Vielleicht erweisen sich zahlreiche numerische Übereinstimmungen als Zufall.


Doch bereits die mathematische Frage ist bemerkenswert:


Wie viel von dem formalen Vokabular der Quantenphysik – Besetzungszahlen, Fermionparität, Antikommutatoren, Symmetrien, Spektren – existiert unabhängig von Physik bereits in der Arithmetik?


Für die quadratfreien Zahlen lautet die Antwort inzwischen überraschend eindeutig: sehr viel.


Und bei \(E_8\) wissen wir seit Viazovska ohnehin, dass eine scheinbar abstrakte mathematische Struktur eine Geometrie von außergewöhnlicher Starrheit besitzt.


Die spannende Forschungsfrage liegt deshalb nicht darin, bekannte Naturkonstanten nachträglich aus \(E_8\) herauszulesen.


Sie lautet vielmehr:


\[
\boxed{
\text{Welche physikähnlichen Strukturen werden durch Mathematik tatsächlich erzwungen?}
}
\]


Erst wenn diese Frage beantwortet ist, darf die zweite gestellt werden:


\[
\boxed{
\text{Hat die Natur diese Struktur gewählt?}
}
\]


Zwischen beiden Fragen liegt die wichtigste Grenze des gesamten Projekts.


Und möglicherweise beginnt genau dort die interessante neue Physik.

\end{document}
